\hypertarget{data-persistence-object-serialization}{%
\chapter{Data Persistence \& Object
Serialization}\label{data-persistence-object-serialization}}

Data persistence allows Python objects to survive the termination of the
process. This involves serialization (converting an object to a byte
stream) and storage.

\hypertarget{pickle-the-virtual-stack-machine}{%
\subsection{\texorpdfstring{36.1 \texttt{pickle}: The Virtual Stack
Machine}{36.1 pickle: The Virtual Stack Machine}}\label{pickle-the-virtual-stack-machine}}

The \passthrough{\lstinline!pickle!} module is Python's native
serialization format. Unlike JSON, it can serialize almost any Python
object (including classes, functions, and complex circular references).

\hypertarget{the-pickle-protocol}{%
\subsubsection{1. The Pickle Protocol}\label{the-pickle-protocol}}

\passthrough{\lstinline!pickle!} doesn't just store data; it stores a
\textbf{program} that, when executed by the pickle virtual machine,
reconstructs the object. * \textbf{Opcodes}: A pickle stream consists of
a series of opcodes (e.g., \passthrough{\lstinline!PROTO!},
\passthrough{\lstinline!EMPTY\_DICT!},
\passthrough{\lstinline!SETITEM!}, \passthrough{\lstinline!STOP!}). *
\textbf{The Stack}: The pickle VM uses a stack to build objects. For
example, to create a list, it might push several items and then call an
opcode that pops them into a new list object.

\hypertarget{security-warning-__reduce__}{%
\subsubsection{\texorpdfstring{2. Security Warning:
\texttt{\_\_reduce\_\_}}{2. Security Warning: \_\_reduce\_\_}}\label{security-warning-__reduce__}}

When an object is unpickled, the VM may execute arbitrary code. The
\passthrough{\lstinline!\_\_reduce\_\_!} method allows an object to
define exactly how it should be reconstructed, which can be exploited to
execute shell commands. \textbf{Never unpickle data from an untrusted
source.}

\hypertarget{json-the-universal-exchange-format}{%
\subsection{\texorpdfstring{36.2 \texttt{json}: The Universal Exchange
Format}{36.2 json: The Universal Exchange Format}}\label{json-the-universal-exchange-format}}

The \passthrough{\lstinline!json!} module provides a standard way to
serialize basic Python types (dicts, lists, strings, numbers) into a
format readable by almost any language.

\hypertarget{c-optimization-_json}{%
\subsubsection{\texorpdfstring{1. C Optimization:
\texttt{\_json}}{1. C Optimization: \_json}}\label{c-optimization-_json}}

In CPython, the \passthrough{\lstinline!json!} module is backed by a C
extension (\passthrough{\lstinline!\_json.c!}). * \textbf{Encoding}:
Iterates through Python objects and builds a C string buffer. *
\textbf{Decoding}: Uses a fast scan-based parser to identify JSON tokens
and convert them to Python objects.

\hypertarget{limitations}{%
\subsubsection{2. Limitations}\label{limitations}}

\passthrough{\lstinline!json!} cannot handle complex Python objects,
circular references, or non-string keys in dictionaries. For these,
custom \passthrough{\lstinline!JSONEncoder!} and
\passthrough{\lstinline!JSONDecoder!} subclasses are required.

\hypertarget{sqlite3-the-embedded-database}{%
\subsection{\texorpdfstring{36.3 \texttt{sqlite3}: The Embedded
Database}{36.3 sqlite3: The Embedded Database}}\label{sqlite3-the-embedded-database}}

Python comes with a complete SQL database engine: SQLite.

\hypertarget{the-c-extension-architecture}{%
\subsubsection{1. The C Extension
Architecture}\label{the-c-extension-architecture}}

The \passthrough{\lstinline!sqlite3!} module is a wrapper around the
SQLite C library. * \textbf{Connection and Cursor}: These are C-level
objects that manage the database file and the result set pointers. *
\textbf{GIL Management}: The \passthrough{\lstinline!sqlite3!} module
releases the GIL during long-running SQL queries, allowing other Python
threads to run while the database is processing I/O or complex joins.

\hypertarget{type-mapping}{%
\subsubsection{2. Type Mapping}\label{type-mapping}}

The module automatically maps Python types to SQL types (e.g.,
\passthrough{\lstinline!int!} to \passthrough{\lstinline!INTEGER!},
\passthrough{\lstinline!str!} to \passthrough{\lstinline!TEXT!}). You
can register custom adapters and converters to handle complex types like
\passthrough{\lstinline!datetime!} or even
\passthrough{\lstinline!pickle!} objects.

\hypertarget{dbm-and-shelve-simple-key-value-stores}{%
\subsection{\texorpdfstring{36.4 \texttt{dbm} and \texttt{shelve}:
Simple Key-Value
Stores}{36.4 dbm and shelve: Simple Key-Value Stores}}\label{dbm-and-shelve-simple-key-value-stores}}

\begin{itemize}
\tightlist
\item
  \textbf{\passthrough{\lstinline!dbm!}}: Provides an interface to Unix
  ``database manager'' libraries (like GDBM or Berkeley DB). It stores
  string keys and values in a disk-based hash table.
\item
  \textbf{\passthrough{\lstinline!shelve!}}: A wrapper around
  \passthrough{\lstinline!dbm!} that uses
  \passthrough{\lstinline!pickle!} to serialize the values. This allows
  you to treat a disk file as a persistent Python dictionary.
\end{itemize}

\begin{center}\rule{0.5\linewidth}{0.5pt}\end{center}

\begin{center}\rule{0.5\linewidth}{0.5pt}\end{center}
